\documentclass{article}

\usepackage[T1]{fontenc}
\usepackage{lmodern}
\usepackage{amsmath}
\usepackage{amssymb}
\usepackage{amsthm}
\usepackage{mathtools}
\usepackage[hidelinks]{hyperref}

\newtheorem*{mainproposition}{Proposition}
\newtheorem{lemma}{Lemma}

\DeclareMathOperator{\wt}{wt}

\title{A local polynomial strip bound}
\author{}
\date{Self-contained hand proof, 11 November 2026}

\begin{document}

\maketitle

\section{Main proposition}\label{sec:main}

\begin{mainproposition}[weight at least fifteen; independent review\slash formalization pending]
  Let \((a_i,b_i)_{i=0}^{18}\) be polynomial pairs over \(\mathbb{F}_2\), with each nonzero
  polynomial of degree at most 18. Assume either of these two systems, with all equalities in the
  indicated polynomial ring:

  \textbf{Forward system F}, in \(\mathbb{F}_2[y]\):
  \begin{equation}\label{eq:forward}
    a_i = (y + y^2) a_{i-2} + b_{i-1} + b_{i-2} + y^2 b_i \quad (2 \le i \le 18).
  \end{equation}

  \textbf{Backward system B}, in \(\mathbb{F}_2[x]\):
  \begin{equation}\label{eq:backward}
    (x + x^2) b_i = a_i + x^2 a_{i-1} + x^2 a_{i-2} + b_{i-2} \quad (2 \le i \le 18).
  \end{equation}

  For a polynomial \(p\), let \(\wt(p)\) count its nonzero coefficients, including
  \(\wt(0) = 0\). Define
  \begin{equation}
    w_i = \wt(a_i) + \wt(b_i), \quad W = \sum_{i=0}^{18} w_i.
  \end{equation}
  Suppose also that
  \begin{equation}\label{eq:adjacent}
    w_i + w_{i+1} \ge 1 \quad (0 \le i \le 17).
  \end{equation}
  Then
  \begin{equation}\label{eq:main-bound}
    W \ge 15.
  \end{equation}
\end{mainproposition}

\textbf{Status and scope.} All lemmas and conclusions in this draft are PROPOSITION, with the named
gap independent review\slash formalization pending. No result is graded THEOREM. The proof below is a
hand argument from the displayed hypotheses alone. It has no charge, code-model,
matrix-recognition, or stabilizer hypothesis and uses no external result about such objects. There
is no enumeration, optimization certificate, or computer proof. In particular the two recurrences
receive separate algebraic proofs.

All polynomial equalities are exact: no quotient relation, periodic identification, or truncation
of products is permitted. The proof does not need the degree cap beyond the statement; it is valid
for arbitrary finite polynomial slices satisfying the same equations and index ranges.

\textbf{Proof organization.} Lemmas~\ref{lem:branch}--\ref{lem:four} establish the local costs.
Section~\ref{sec:counting} combines them by counting gaps between zero slices and completes the
proof.

\section{Polynomial facts and a branch between zeros}\label{sec:polynomial}

For the local calculations only, rename the indeterminate in either system as \(T\), and put
\(t = T + T^2\). This renames a variable, not a recurrence. A slice is empty if \(w_i = 0\), a
singleton if \(w_i = 1\), and heavy if \(w_i \ge 2\). A singleton pair has exactly one nonzero
block, which is a monomial. In particular \(u + T^2 v\) and \(T^2 u + v\) each have weight one
whenever \((u,v)\) is a singleton.

If \(\rho(p)\) counts maximal consecutive runs of exponents in the support of \(p\), with
\(\rho(0) = 0\), then
\begin{equation}\label{eq:runs}
  \wt(tp) = \wt((1 + T)p) = 2\rho(p).
\end{equation}
Indeed a run from exponent \(c\) to \(d\) contributes exactly \(T^c + T^{d+1}\) to \((1 + T)p\).
Endpoints from distinct runs are distinct. A monomial shift does not change weight. Thus
\(\wt(tp) \ge 2\) for \(p \ne 0\). We also use \(\wt(u + v) \le \wt(u) + \wt(v)\) and the fact
that \(p(1)\) is the parity of \(\wt(p)\). Both follow coefficient by coefficient over
\(\mathbb{F}_2\). The polynomial ring is a domain, so nonzero factors can be cancelled.

\begin{lemma}[directional branch]\label{lem:branch}
  Suppose \(1 \le r \le 17\) and slices \(r - 1, r + 1\) are empty. In system F, \(b_r = 0\); if
  \(r \le 16\) and slice \(r\) is nonempty, then \(w_{r+2} \ge 2\). In system B, \(a_r = 0\); if
  \(r \ge 2\) and slice \(r\) is nonempty, then \(w_{r-2} \ge 2\).
\end{lemma}

\begin{proof}[Forward proof]
  Applying equation~\eqref{eq:forward} at \(r + 1\) gives \(0 = b_r\). At \(r + 2\), when this
  index is at most 18, it gives
  \begin{equation}
    a_{r+2} + T^2 b_{r+2} = t a_r.
  \end{equation}
  The sum of the two physical block weights is at least the weight of their sum, hence at least
  \(2\rho(a_r) \ge 2\) when \(a_r \ne 0\).
\end{proof}

\begin{proof}[Backward proof]
  Applying equation~\eqref{eq:backward} at \(r + 1\) gives \(0 = T^2 a_r\), so \(a_r = 0\).
  Applying equation~\eqref{eq:backward} at \(r\), when \(r \ge 2\), gives
  \begin{equation}
    T^2 a_{r-2} + b_{r-2} = t b_r.
  \end{equation}
  The same coefficient-weight inequality gives \(w_{r-2} \ge 2\rho(b_r) \ge 2\). This proves both
  assertions.
\end{proof}

There is exactly one possible exceptional middle index per direction: \(r = 17\) in F would
require equation 19, and \(r = 1\) in B would require equation 1. Neither equation is among the
hypotheses. No branch cost is asserted at these exceptions.

\section{Local costs inside gaps}\label{sec:gaps}

\begin{lemma}[two occupied slices between zeros]\label{lem:two}
  If \(0 \le k \le 15\), slices \(k, k + 3\) are empty, and the intermediate slices are not both
  empty, then
  \begin{equation}
    w_{k+1} + w_{k+2} \ge 4.
  \end{equation}
\end{lemma}

\begin{proof}[Forward proof]
  Write the intermediate slices as \((p,q), (u,v)\). Equations at \(k + 2, k + 3\) are
  \begin{equation}
    u = q + T^2 v, \quad tp + q + v = 0.
  \end{equation}
  Evaluation at \(T = 1\) gives \(u(1) = 0\). If \(u \ne 0\), then \(\wt(u) \ge 2\). Also
  \(\wt(p) + \wt(q) + \wt(v) \ge 2\): weight zero forces \(u = 0\), and weight one leaves a single
  nonzero term in \(tp + q + v = 0\). Thus the total is at least four. If \(u = 0\), then
  \(q = T^2 v\) and \(tp = (1 + T^2)v\). Cancelling \(1 + T\) yields \(Tp = (1 + T)v\). Evaluation
  at zero shows that \(v = Th\) for a polynomial \(h\); then \(p = (1 + T)h\) and \(q = T^3 h\).
  Nonemptiness gives \(h \ne 0\), and total weight is \(2\wt(h) + 2\rho(h) \ge 4\).
\end{proof}

\begin{proof}[Backward proof]
  The equations for the same named slices are instead
  \begin{equation}
    tv = u + T^2 p, \quad q = T^2 (u + p).
  \end{equation}
  Evaluation at one gives \(q(1) = 0\). If \(q \ne 0\), its weight is at least two. The remaining
  weight \(\wt(p) + \wt(u) + \wt(v)\) is at least two, since weight zero forces \(q = 0\) and
  weight one cannot satisfy \(tv + u + T^2 p = 0\). If \(q = 0\), then \(u = p\) and
  \(tv = (1 + T^2)p\). Cancelling \(1 + T\) gives \(Tv = (1 + T)p\). Hence \(p = Th\), \(u = Th\),
  \(v = (1 + T)h\) for \(h \ne 0\). The total is again \(2\wt(h) + 2\rho(h) \ge 4\).
\end{proof}

\begin{lemma}[three occupied slices between zeros]\label{lem:three}
  If \(0 \le k \le 14\), slices \(k, k + 4\) are empty, and all three intermediate slices are
  occupied, then
  \begin{equation}
    w_{k+1} + w_{k+2} + w_{k+3} \ge 4.
  \end{equation}
\end{lemma}

\begin{proof}[Forward proof]
  A failure would make all three pairs \((p,q), (u,v), (h,l)\) singletons. The equations at
  \(k + 2, k + 3, k + 4\) are
  \begin{equation}
    u = q + T^2 v, \quad h = tp + v + q + T^2 l, \quad 0 = tu + l + v.
  \end{equation}
  If \(q = 0\), the first equation forces both \(u, v\) to vanish or both to be nonzero,
  contradicting the second singleton. Thus \(q\) is a monomial and \(p = 0\). If \(u \ne 0\), then
  \(v = 0\) and the last equation gives \(l = tu\) of weight two, contradicting the third
  singleton. Consequently \(u = 0\), \(v\) is a monomial, and \(l = v\), so \(h = 0\). The middle
  equation now gives \(q = (1 + T^2)v\), of weight two, a contradiction.
\end{proof}

\begin{proof}[Backward proof]
  The equations for the three singleton slices are
  \begin{equation}
    tv = u + T^2 p, \quad tl = h + T^2 u + T^2 p + q, \quad v = T^2 (h + u).
  \end{equation}
  If \(v \ne 0\), it is a monomial and \(u = 0\), so the first equation forces \(T^2 p = tv\) of
  weight two, impossible in the first singleton. Hence \(v = 0\) and \(u\) is a monomial. The last
  equation gives \(h = u\), so \(l = 0\). The first gives \(T^2 p = u\), so \(q = 0\). The middle
  equation is then \(0 = u + T^2 u + u = T^2 u\), a contradiction.
\end{proof}

\begin{lemma}[four occupied slices next to a zero]\label{lem:four}
  In system F, an empty slice \(k\), with \(0 \le k \le 14\), followed by four occupied slices
  satisfies
  \begin{equation}
    \sum_{i=k+1}^{k+4} w_i \ge 5.
  \end{equation}
  In system B, an empty slice \(R\), with \(4 \le R \le 18\), preceded by four occupied slices
  satisfies
  \begin{equation}
    \sum_{i=R-4}^{R-1} w_i \ge 5.
  \end{equation}
  Neither assertion requires another empty slice at the other end.
\end{lemma}

\begin{proof}[Forward proof]
  If all four slices \((p,q), (u,v), (h,l), (m,n)\) are singletons in increasing index order, the
  equations at \(k + 2, k + 3, k + 4\) are
  \begin{equation}
    u + T^2 v = q, \quad h + T^2 l = tp + v + q, \quad m + T^2 n = tu + l + v.
  \end{equation}
  The first left side is a monomial, so \(q\) is a monomial and \(p = 0\). If the second slice
  were B-only, then \(u = 0, q = T^2 v\), and the second equation would give
  \(h + T^2 l = (1 + T^2)v\) of weight two, a contradiction. Thus the second slice is A-only, with
  \(u = q, v = 0\). The third equation to solve for a slice is now \(h + T^2 l = q\). If the third
  slice is A-only, then \(h = q, l = 0\), and the fourth-slice equation has right side \(tq\), of
  weight two. If it is B-only, then \(h = 0, q = T^2 l\), and that right side is
  \(tq + l = (1 + T^3 + T^4)l\), of weight three. Both contradict the fourth singleton, whose left
  side \(m + T^2 n\) is a monomial.
\end{proof}

\begin{proof}[Backward proof]
  Name singleton slices \((p,q), (u,v), (h,l), (m,n)\) in decreasing order at
  \(R - 1, R - 2, R - 3, R - 4\). The equations at \(R, R - 1, R - 2\) are
  \begin{equation}
    T^2 u + v = T^2 p, \quad T^2 h + l = tq + p + T^2 u, \quad T^2 m + n = tv + u + T^2 h.
  \end{equation}
  The first left side is a monomial, so \(p\) is a monomial and \(q = 0\). If the second slice
  were A-only, then \(u = p, v = 0\), and the second equation would give
  \(T^2 h + l = (1 + T^2)p\) of weight two. Thus the second slice is B-only, with
  \(u = 0, v = T^2 p\). The next equation is \(T^2 h + l = p\). If the third slice is B-only, then
  \(h = 0, l = p\), and the last right side is \(t T^2 p\), of weight two. If it is A-only, then
  \(l = 0, T^2 h = p\), and the last right side is \((1 + T^3 + T^4)p\), of weight three. Neither
  is the monomial \(T^2 m + n\) required by the fourth singleton.
\end{proof}

In a gap of length five between successive zero slices, all four interior slices are occupied.
Lemma~\ref{lem:four} applies from the left endpoint in F and from the right endpoint in B. These
endpoints satisfy \(k \le 13\) and \(R \ge 5\), respectively. Thus such a gap has interior weight
at least five. All monomials counted in the proofs have distinct exponents; none may be discarded
if an exponent exceeds 18. Such an identity would violate the hypotheses rather than supply a
further cancellation.

\section{Gap counting and completion of the main proof}\label{sec:counting}

Let \(z\) be the number of zero slices, and let
\begin{equation}
  E = \sum_{i : w_i > 0} (w_i - 1) = W - (19 - z).
\end{equation}
This is a nonnegative integer. If \(z \le 1\), then \(W \ge 18\), so assume \(z \ge 2\). List
zero positions as \(q_1 < \ldots < q_z\). By \eqref{eq:adjacent}, every gap \(q_{j+1} - q_j\) is
at least two and every interior slice of a successive-zero gap is occupied. Let \(G, H, K, L\)
count gaps of lengths two, three, four, and five respectively.

\textbf{Branch budget.} The middle slice \(r\) of a length-two gap is nonempty. By
Lemma~\ref{lem:branch} it forces a heavy slice at \(r + 2\) in F or \(r - 2\) in B, with at most
one exception (respectively \(r = 17\) or \(r = 1\)). Distinct middle indices have distinct
targets. Each target contributes at least one unit to \(E\), irrespective of whether the middle
slice is itself heavy. Therefore
\begin{equation}\label{eq:branch-budget}
  G \le E + 1.
\end{equation}

\textbf{Interior budget.} A length-three gap costs at least two units of excess by
Lemma~\ref{lem:two}. A length-four gap costs at least one by Lemma~\ref{lem:three}, and a
length-five gap at least one by Lemma~\ref{lem:four}. Interiors of distinct successive-zero gaps
are disjoint. Hence
\begin{equation}\label{eq:interior-budget}
  2H + K + L \le E.
\end{equation}

\textbf{Span budget.} The total span \(q_z - q_1\) is at most 18. All uncounted gaps have length
at least six, so
\begin{equation}\label{eq:span-budget}
  \begin{aligned}
    18 &\ge 2G + 3H + 4K + 5L + 6(z - 1 - G - H - K - L), \\
    4G + 3H + 2K + L &\ge 6z - 24.
  \end{aligned}
\end{equation}

Since \(3H + 2K + L \le 2(2H + K + L) \le 2E\), combining \eqref{eq:branch-budget},
\eqref{eq:interior-budget}, and \eqref{eq:span-budget} yields
\begin{equation}
  6z - 24 \le 4G + 3H + 2K + L \le 4(E + 1) + 2E = 6E + 4.
\end{equation}
The two excess budgets may overlap: the coefficient six explicitly allows this. No disjointness
between branch targets and gap interiors was used. Consequently
\begin{equation}
  6E \ge 6z - 28, \quad 6W = 114 - 6z + 6E \ge 86.
\end{equation}
Integer weight now gives \(W \ge \lceil 86/6 \rceil = 15\), completing the main proof. \qed

\textbf{Review boundary.} This draft contains only the local statement and its proof. It makes no
claim of a bound twenty and no assertion connecting these hypotheses to any external model. All
claims remain PROPOSITION pending independent review and formalization.

\end{document}
